\textit{The source also provides this cloudy-weather observation task.}
You have 30 minutes. The telescope and finder tube are already aligned. Do not change the alignment.
Both right-ascension and declination dials are calibrated. Write answers on the separate answer sheet.
Answers/calculations on the question paper will not be evaluated.
The north--south line is parallel to the balcony railing, with north on your left.
\ioaafig{2023_OBS-TEL_Q16-f1.pdf}
Point A is a summit behind the first row of mountains. Point B is the tallest summit in the first row,
and point C is a square building at the right end of the mountain.
\begin{description}
\item[(A)] Point the telescope to the three arrowed points and use its dials to determine their equatorial coordinates. \hfill\mbox{[15 points]}
\item[(B)] Find the angular distance between points A and C. \hfill\mbox{[10 points]}
\item[(C)] Call the examiner and demonstrate how to ``scan'' the celestial equator with the telescope,
passing through all equatorial constellations. Give the examiner your answer sheet to record the result. \hfill\mbox{[10 points]}
\item[(D)] Point the telescope towards equatorial coordinates $13^{\mathrm h}20^{\mathrm m}$ and $2^\circ$ S.
Ask the examiner to confirm successful targeting. \hfill\mbox{[10 points]}
\end{description}
